% Общая преамбула лекций курса. Подключается после \documentclass:
%
%     \documentclass[aspectratio=169,12pt]{beamer}
%     % Общая преамбула лекций курса. Подключается после \documentclass:
%
%     \documentclass[aspectratio=169,12pt]{beamer}
%     % Общая преамбула лекций курса. Подключается после \documentclass:
%
%     \documentclass[aspectratio=169,12pt]{beamer}
%     % Общая преамбула лекций курса. Подключается после \documentclass:
%
%     \documentclass[aspectratio=169,12pt]{beamer}
%     \input{mlg_preamble}
%
% Здесь только оформление и макросы — ни \title, ни текста слайдов, чтобы
% пять презентаций курса выглядели одинаково и правились в одном месте.
%
\usepackage{fontspec}
\usepackage{polyglossia}
\setmainlanguage{russian}
\setotherlanguage{english}
\setmainfont{Calibri}
\setsansfont{Calibri}
\setmonofont{Consolas}[Scale=0.92]

\usepackage{graphicx}
\usepackage{booktabs}
\usepackage{array}

% ---------------------------------------------------------------- оформление
\definecolor{ink}{HTML}{1B2433}
\definecolor{navy}{HTML}{2B4C6F}
\definecolor{brick}{HTML}{B03A2B}
\definecolor{moss}{HTML}{3E7D5A}
\definecolor{plum}{HTML}{6B5B95}
\definecolor{ash}{HTML}{8B9299}
\definecolor{paper}{HTML}{ECEFF3}

\usetheme{default}
\usecolortheme{default}
\setbeamercolor{structure}{fg=navy}
\setbeamercolor{frametitle}{fg=navy}
\setbeamercolor{normal text}{fg=ink}
\setbeamercolor{itemize item}{fg=navy}
\setbeamercolor{itemize subitem}{fg=ash}
\setbeamerfont{frametitle}{size=\Large,series=\bfseries}
\setbeamerfont{framesubtitle}{size=\small,series=\mdseries}
\setbeamercolor{framesubtitle}{fg=ash}
\setbeamertemplate{navigation symbols}{}
\setbeamertemplate{itemize item}{\raisebox{0.15ex}{\tiny$\bullet$}}
\setbeamertemplate{itemize subitem}{\raisebox{0.15ex}{\tiny--}}
\setbeamersize{text margin left=1.1cm, text margin right=1.1cm}

\setbeamertemplate{footline}{%
  \hbox{\begin{beamercolorbox}[wd=\paperwidth,ht=2.6ex,dp=1.4ex,leftskip=1.1cm,%
        rightskip=1.1cm]{}%
    {\color{ash}\scriptsize Основы машинного обучения \quad·\quad
     Цифровая лингвистика и локализация}\hfill
    {\color{ash}\scriptsize\insertframenumber}%
  \end{beamercolorbox}}}

% Врезка: связь слайда с лабораторной или с работой переводчика.
% Ширина считается точно: colorbox добавляет \fboxsep с каждой стороны,
% поэтому parbox должен быть на 2\fboxsep уже текстового поля — иначе
% на каждом слайде вылезает overfull hbox.
\newcommand{\врезка}[1]{%
  \par\vspace{0.5ex}\noindent
  \colorbox{paper}{%
    \parbox{\dimexpr\textwidth-2\fboxsep\relax}{%
      \vspace{0.7ex}\small #1\vspace{0.7ex}}}\par}

% Колонка таблицы заданной доли ширины, без растягивания пробелов.
% Объявлять её нужно здесь: внутри frame beamer пересканирует тело кадра,
% и #1 в определении вызывает «Illegal parameter number».
\newcolumntype{Л}[1]{>{\raggedright\arraybackslash}p{#1\textwidth}}

% Ссылка на интерактивную страницу к слайду. Playground'ы считают те же
% числа на том же корпусе, что и лабораторные, поэтому на них можно
% ссылаться прямо с лекции, не оговаривая расхождений.
\hypersetup{colorlinks=true,urlcolor=moss,linkcolor=moss}
\newcommand{\playground}[3]{%
  \par\vspace{0.4ex}%
  {\footnotesize\color{ash}$\rightarrow$~\href{#1}{\bfseries #2} — #3\par}}

\newcommand{\кейс}[1]{{\color{brick}\small\textbf{#1}}}
\newcommand{\числ}[1]{{\color{navy}\bfseries #1}}

%
% Здесь только оформление и макросы — ни \title, ни текста слайдов, чтобы
% пять презентаций курса выглядели одинаково и правились в одном месте.
%
\usepackage{fontspec}
\usepackage{polyglossia}
\setmainlanguage{russian}
\setotherlanguage{english}
\setmainfont{Calibri}
\setsansfont{Calibri}
\setmonofont{Consolas}[Scale=0.92]

\usepackage{graphicx}
\usepackage{booktabs}
\usepackage{array}

% ---------------------------------------------------------------- оформление
\definecolor{ink}{HTML}{1B2433}
\definecolor{navy}{HTML}{2B4C6F}
\definecolor{brick}{HTML}{B03A2B}
\definecolor{moss}{HTML}{3E7D5A}
\definecolor{plum}{HTML}{6B5B95}
\definecolor{ash}{HTML}{8B9299}
\definecolor{paper}{HTML}{ECEFF3}

\usetheme{default}
\usecolortheme{default}
\setbeamercolor{structure}{fg=navy}
\setbeamercolor{frametitle}{fg=navy}
\setbeamercolor{normal text}{fg=ink}
\setbeamercolor{itemize item}{fg=navy}
\setbeamercolor{itemize subitem}{fg=ash}
\setbeamerfont{frametitle}{size=\Large,series=\bfseries}
\setbeamerfont{framesubtitle}{size=\small,series=\mdseries}
\setbeamercolor{framesubtitle}{fg=ash}
\setbeamertemplate{navigation symbols}{}
\setbeamertemplate{itemize item}{\raisebox{0.15ex}{\tiny$\bullet$}}
\setbeamertemplate{itemize subitem}{\raisebox{0.15ex}{\tiny--}}
\setbeamersize{text margin left=1.1cm, text margin right=1.1cm}

\setbeamertemplate{footline}{%
  \hbox{\begin{beamercolorbox}[wd=\paperwidth,ht=2.6ex,dp=1.4ex,leftskip=1.1cm,%
        rightskip=1.1cm]{}%
    {\color{ash}\scriptsize Основы машинного обучения \quad·\quad
     Цифровая лингвистика и локализация}\hfill
    {\color{ash}\scriptsize\insertframenumber}%
  \end{beamercolorbox}}}

% Врезка: связь слайда с лабораторной или с работой переводчика.
% Ширина считается точно: colorbox добавляет \fboxsep с каждой стороны,
% поэтому parbox должен быть на 2\fboxsep уже текстового поля — иначе
% на каждом слайде вылезает overfull hbox.
\newcommand{\врезка}[1]{%
  \par\vspace{0.5ex}\noindent
  \colorbox{paper}{%
    \parbox{\dimexpr\textwidth-2\fboxsep\relax}{%
      \vspace{0.7ex}\small #1\vspace{0.7ex}}}\par}

% Колонка таблицы заданной доли ширины, без растягивания пробелов.
% Объявлять её нужно здесь: внутри frame beamer пересканирует тело кадра,
% и #1 в определении вызывает «Illegal parameter number».
\newcolumntype{Л}[1]{>{\raggedright\arraybackslash}p{#1\textwidth}}

% Ссылка на интерактивную страницу к слайду. Playground'ы считают те же
% числа на том же корпусе, что и лабораторные, поэтому на них можно
% ссылаться прямо с лекции, не оговаривая расхождений.
\hypersetup{colorlinks=true,urlcolor=moss,linkcolor=moss}
\newcommand{\playground}[3]{%
  \par\vspace{0.4ex}%
  {\footnotesize\color{ash}$\rightarrow$~\href{#1}{\bfseries #2} — #3\par}}

\newcommand{\кейс}[1]{{\color{brick}\small\textbf{#1}}}
\newcommand{\числ}[1]{{\color{navy}\bfseries #1}}

%
% Здесь только оформление и макросы — ни \title, ни текста слайдов, чтобы
% пять презентаций курса выглядели одинаково и правились в одном месте.
%
\usepackage{fontspec}
\usepackage{polyglossia}
\setmainlanguage{russian}
\setotherlanguage{english}
\setmainfont{Calibri}
\setsansfont{Calibri}
\setmonofont{Consolas}[Scale=0.92]

\usepackage{graphicx}
\usepackage{booktabs}
\usepackage{array}

% ---------------------------------------------------------------- оформление
\definecolor{ink}{HTML}{1B2433}
\definecolor{navy}{HTML}{2B4C6F}
\definecolor{brick}{HTML}{B03A2B}
\definecolor{moss}{HTML}{3E7D5A}
\definecolor{plum}{HTML}{6B5B95}
\definecolor{ash}{HTML}{8B9299}
\definecolor{paper}{HTML}{ECEFF3}

\usetheme{default}
\usecolortheme{default}
\setbeamercolor{structure}{fg=navy}
\setbeamercolor{frametitle}{fg=navy}
\setbeamercolor{normal text}{fg=ink}
\setbeamercolor{itemize item}{fg=navy}
\setbeamercolor{itemize subitem}{fg=ash}
\setbeamerfont{frametitle}{size=\Large,series=\bfseries}
\setbeamerfont{framesubtitle}{size=\small,series=\mdseries}
\setbeamercolor{framesubtitle}{fg=ash}
\setbeamertemplate{navigation symbols}{}
\setbeamertemplate{itemize item}{\raisebox{0.15ex}{\tiny$\bullet$}}
\setbeamertemplate{itemize subitem}{\raisebox{0.15ex}{\tiny--}}
\setbeamersize{text margin left=1.1cm, text margin right=1.1cm}

\setbeamertemplate{footline}{%
  \hbox{\begin{beamercolorbox}[wd=\paperwidth,ht=2.6ex,dp=1.4ex,leftskip=1.1cm,%
        rightskip=1.1cm]{}%
    {\color{ash}\scriptsize Основы машинного обучения \quad·\quad
     Цифровая лингвистика и локализация}\hfill
    {\color{ash}\scriptsize\insertframenumber}%
  \end{beamercolorbox}}}

% Врезка: связь слайда с лабораторной или с работой переводчика.
% Ширина считается точно: colorbox добавляет \fboxsep с каждой стороны,
% поэтому parbox должен быть на 2\fboxsep уже текстового поля — иначе
% на каждом слайде вылезает overfull hbox.
\newcommand{\врезка}[1]{%
  \par\vspace{0.5ex}\noindent
  \colorbox{paper}{%
    \parbox{\dimexpr\textwidth-2\fboxsep\relax}{%
      \vspace{0.7ex}\small #1\vspace{0.7ex}}}\par}

% Колонка таблицы заданной доли ширины, без растягивания пробелов.
% Объявлять её нужно здесь: внутри frame beamer пересканирует тело кадра,
% и #1 в определении вызывает «Illegal parameter number».
\newcolumntype{Л}[1]{>{\raggedright\arraybackslash}p{#1\textwidth}}

% Ссылка на интерактивную страницу к слайду. Playground'ы считают те же
% числа на том же корпусе, что и лабораторные, поэтому на них можно
% ссылаться прямо с лекции, не оговаривая расхождений.
\hypersetup{colorlinks=true,urlcolor=moss,linkcolor=moss}
\newcommand{\playground}[3]{%
  \par\vspace{0.4ex}%
  {\footnotesize\color{ash}$\rightarrow$~\href{#1}{\bfseries #2} — #3\par}}

\newcommand{\кейс}[1]{{\color{brick}\small\textbf{#1}}}
\newcommand{\числ}[1]{{\color{navy}\bfseries #1}}

%
% Здесь только оформление и макросы — ни \title, ни текста слайдов, чтобы
% пять презентаций курса выглядели одинаково и правились в одном месте.
%
\usepackage{fontspec}
\usepackage{polyglossia}
\setmainlanguage{russian}
\setotherlanguage{english}
\setmainfont{Calibri}
\setsansfont{Calibri}
\setmonofont{Consolas}[Scale=0.92]

\usepackage{graphicx}
\usepackage{booktabs}
\usepackage{array}

% ---------------------------------------------------------------- оформление
\definecolor{ink}{HTML}{1B2433}
\definecolor{navy}{HTML}{2B4C6F}
\definecolor{brick}{HTML}{B03A2B}
\definecolor{moss}{HTML}{3E7D5A}
\definecolor{plum}{HTML}{6B5B95}
\definecolor{ash}{HTML}{8B9299}
\definecolor{paper}{HTML}{ECEFF3}

\usetheme{default}
\usecolortheme{default}
\setbeamercolor{structure}{fg=navy}
\setbeamercolor{frametitle}{fg=navy}
\setbeamercolor{normal text}{fg=ink}
\setbeamercolor{itemize item}{fg=navy}
\setbeamercolor{itemize subitem}{fg=ash}
\setbeamerfont{frametitle}{size=\Large,series=\bfseries}
\setbeamerfont{framesubtitle}{size=\small,series=\mdseries}
\setbeamercolor{framesubtitle}{fg=ash}
\setbeamertemplate{navigation symbols}{}
\setbeamertemplate{itemize item}{\raisebox{0.15ex}{\tiny$\bullet$}}
\setbeamertemplate{itemize subitem}{\raisebox{0.15ex}{\tiny--}}
\setbeamersize{text margin left=1.1cm, text margin right=1.1cm}

\setbeamertemplate{footline}{%
  \hbox{\begin{beamercolorbox}[wd=\paperwidth,ht=2.6ex,dp=1.4ex,leftskip=1.1cm,%
        rightskip=1.1cm]{}%
    {\color{ash}\scriptsize Основы машинного обучения \quad·\quad
     Цифровая лингвистика и локализация}\hfill
    {\color{ash}\scriptsize\insertframenumber}%
  \end{beamercolorbox}}}

% Врезка: связь слайда с лабораторной или с работой переводчика.
% Ширина считается точно: colorbox добавляет \fboxsep с каждой стороны,
% поэтому parbox должен быть на 2\fboxsep уже текстового поля — иначе
% на каждом слайде вылезает overfull hbox.
\newcommand{\врезка}[1]{%
  \par\vspace{0.5ex}\noindent
  \colorbox{paper}{%
    \parbox{\dimexpr\textwidth-2\fboxsep\relax}{%
      \vspace{0.7ex}\small #1\vspace{0.7ex}}}\par}

% Колонка таблицы заданной доли ширины, без растягивания пробелов.
% Объявлять её нужно здесь: внутри frame beamer пересканирует тело кадра,
% и #1 в определении вызывает «Illegal parameter number».
\newcolumntype{Л}[1]{>{\raggedright\arraybackslash}p{#1\textwidth}}

% Ссылка на интерактивную страницу к слайду. Playground'ы считают те же
% числа на том же корпусе, что и лабораторные, поэтому на них можно
% ссылаться прямо с лекции, не оговаривая расхождений.
\hypersetup{colorlinks=true,urlcolor=moss,linkcolor=moss}
\newcommand{\playground}[3]{%
  \par\vspace{0.4ex}%
  {\footnotesize\color{ash}$\rightarrow$~\href{#1}{\bfseries #2} — #3\par}}

\newcommand{\кейс}[1]{{\color{brick}\small\textbf{#1}}}
\newcommand{\числ}[1]{{\color{navy}\bfseries #1}}
